\section{Sequential schedules and proof-graph compilation}\label{sec:schedule}

\subsection{From DAGs to schedules}

A finite proof graph can be topologically sorted.  Rather than formalizing arbitrary graph theory in
the logical kernel, \TAPIR{} uses a sequential schedule whose order is produced by an untrusted host
pass and verified by its type.

Fix a residual context \(R\).  A schedule records a list \(O\) of already produced formulas, in
reverse chronological order.  It is generated by:

\begin{align*}
  \mathsf{nil}&:\mathsf{Schedule}(R,[]),\\
  \mathsf{cons}(S,p)&:\mathsf{Schedule}(R,\varphi::O),
\end{align*}
where \(S:\mathsf{Schedule}(R,O)\) and the local plan
\[
  p:\Plan(O++R,\varphi)
\]
may use any earlier output and any residual hypothesis.  If a natural local proof uses only a subset
of \(O++R\), it is weakened to this canonical context.

The reverse order is convenient: the newest output is the head, while all of its dependencies are in
the tail.  A DAG with several roots can be represented by adding a final node that conjoins or
otherwise packages the desired outputs.

\subsection{Compilation theorem}

\begin{theorem}[Schedule compilation]\label{thm:schedule-compile}
Let \(S:\mathsf{Schedule}(R,O)\).  For every \(\psi\in O\), there is a plan
\[
  \mathsf{compile}(S,\psi):\Plan(R,\psi).
\]
Consequently, if \(O=\varphi::O'\), the head has a plan from the residuals alone.
\end{theorem}

\begin{proof}
Induct on \(S\).

For the empty schedule there is no \(\psi\in[]\), so the assertion is vacuous.

Suppose \(S=\mathsf{cons}(S_0,p)\), where
\(S_0:\mathsf{Schedule}(R,O_0)\) and
\(p:\Plan(O_0++R,\varphi)\).  Let \(\psi\in\varphi::O_0\).

If \(\psi=\varphi\), apply the cut constructor to \(p\).  Every hypothesis
\(\chi\in O_0++R\) is discharged as follows.  If \(\chi\in O_0\), the induction hypothesis supplies
\(\Plan(R,\chi)\).  If \(\chi\in R\), use the hypothesis plan
\(\mathsf{prf}(\mathsf{Prf.hyp})\).  Simultaneous cut therefore yields
\(\Plan(R,\varphi)\).

If \(\psi\in O_0\), use the induction hypothesis directly.  These cases exhaust membership in the
cons list.
\end{proof}

The proof is constructive and mirrors a short Lean recursion.  It is the last-mile graph compiler
that prior packets often described semantically but did not connect to the ordinary context calculus.

\begin{corollary}[Acyclic contract graph]\label{cor:dag}
Let a finite directed acyclic graph have nodes labeled by formulas.  Suppose every node \(v\) carries
a local plan deriving its label from the labels of its predecessors and residual context \(R\).  Then
every root label is derivable from \(R\) in \(\PRAAx\).
\end{corollary}

\begin{proof}
Choose any topological order and build a schedule.  At each stage, all predecessors are earlier
outputs.  Weaken the local plan from its actual dependency list to the full list of earlier outputs
plus \(R\).  Apply \cref{thm:schedule-compile} to the root node.
\end{proof}

\subsection{Closing residuals}

\begin{theorem}[Residual closure]\label{thm:residual-close}
If \(R\vdash_{\PRAAx}\varphi\) and for every \(B\in R\) one has
\(\Delta\vdash_{\PRAAx}B\), then \(\Delta\vdash_{\PRAAx}\varphi\).
\end{theorem}

\begin{proof}
This is exactly simultaneous cut, \code{Prf.trans}.  No induction on the mathematical proof graph is
needed because the schedule compiler has already reduced the graph to one ordinary derivation from
\(R\).
\end{proof}

A useful consequence is monotonic development.  A packet may initially expose five residuals.  A
later theorem can discharge two of them, producing a stronger packet with three residuals by a
single cut, without regenerating the internal proof.

\subsection{Exactness of the residual list}

The phrase ``exact residuals'' is syntactic, not a claim of logical independence.  The host compiler
performs a backward slice from the selected root, removes unreachable nodes, and takes the union of
external boundaries actually mentioned by the reachable local plans.  The resulting schedule has
that list as its context.  Therefore:

\begin{proposition}[Syntactic residual exactness]\label{prop:residual-exact}
For a sliced schedule generated by the reference algorithm, every formula in the reported residual
list occurs in at least one reachable local context, and no external formula absent from that list is
used by the compiled plan.
\end{proposition}

\begin{proof}
The first assertion follows from how the union is accumulated during the backward traversal; a
boundary is inserted only when a reachable node declares it.  The second follows by induction on the
schedule: each local plan is typed over earlier outputs plus the reported residual context, and the
compiler introduces no new hypotheses.  This does not prove that a residual is mathematically
indispensable; another proof may avoid it.
\end{proof}

\subsection{Dead-code elimination}

Let \(G\) be a finite graph and \(r\) a root.  Define \(\mathrm{reach}(r)\) as the least set containing
\(r\) and every predecessor of a member.  Delete all other nodes before topological sorting.

\begin{lemma}[Backward slicing preserves the root proof]
If every dependency named by a reachable node is retained, deleting unreachable nodes does not
change the local proof of any reachable node and hence does not change the compiled root plan.
\end{lemma}

\begin{proof}
By induction on a topological order of the reachable subgraph.  A reachable node's local plan uses
only declared predecessors and residuals.  Every declared predecessor is reachable by definition,
so its compiled plan remains available.  Unreachable outputs are never selected in the discharge
function used by cut.
\end{proof}

This elementary optimization matters in large packets: it prevents the final object proof from
containing unrelated certificates or producer assumptions.

\subsection{Cycles and demand closure}

The proof schedule itself is acyclic.  A mathematical construction may have cyclic resource
dependencies---for example, a requested precision at one node increases the requirement at another.
Those cycles are resolved in a separate finite host pass.  On a product of finite chains, iterate
\[
 d_{t+1}=d_t\vee F(d_t)
\]
from the least demand.  If \(F\) is monotone, the sequence stabilizes at the least prefixed point;
each strict iteration increases at least one coordinate, so no more than the sum of coordinate
capacities is needed.  The stabilized numerical demand is then frozen and inserted into the acyclic
proof schedule.

If the closure itself is part of the public theorem, a clocked object trace can verify it.  Usually it
is merely a compile-time resource calculation, in which case Lean checks only the specialized proof
that uses the resulting constants.

\subsection{Host algorithm and verification conditions}

The reference compiler performs:

\begin{enumerate}
\item parse nodes, dependencies, conclusions, and residual declarations;
\item reject duplicate identifiers, missing dependencies, and cycles;
\item backward-slice from the selected root;
\item topologically order the slice;
\item compute the residual union and certificate escape modes;
\item emit a schedule description and, eventually, Lean source that constructs the corresponding
      \(\mathsf{Schedule}\) value.
\end{enumerate}

Only the final Lean value matters for soundness.  The Python implementation in the archive tests the
finite algorithms and produces JSON examples, but it is explicitly not a checker for
\(\Prf\,\PRAAx\).

\subsection{Lean form of the schedule theorem}

The index-safe list formulation avoids formula-code equality and arbitrary graph indices:

\begin{lstlisting}[language=Lean]
inductive Schedule {n : Nat} (R : List (Fm PTm n)) :
    List (Fm PTm n) -> Type
  | nil : Schedule R []
  | cons {outs : List (Fm PTm n)} {phi : Fm PTm n}
      (tail : Schedule R outs)
      (local : Plan (outs ++ R) phi) : Schedule R (phi :: outs)

private def dischargeAppend
    (ih : forall psi, psi in outs -> Plan R psi) :
    forall psi, psi in outs ++ R -> Plan R psi := by
  intro psi h
  rcases List.mem_append.mp h with hOut | hR
  case inl => exact ih psi hOut
  case inr => exact Plan.prf (Prf.hyp hR)

def Schedule.proveAll :
    (s : Schedule R outs) -> forall psi, psi in outs -> Plan R psi
  | .nil, _, h => by simp at h
  | .cons tail local, psi, h => by
      rcases List.mem_cons.mp h with rfl | hTail
      case inl => exact Plan.cut local (dischargeAppend (Schedule.proveAll tail))
      case inr => exact Schedule.proveAll tail psi hTail
\end{lstlisting}

This file is supplied without placeholders.  It still requires elaboration against the pinned
checkout before any verification claim can be made.
